\ifdefined\gkver\else\def\gkver{student}\fi
\documentclass[\gkver]{gaokaozhenti}
\begin{document}

\chapter{2008年广东理科基础}
%% paperType: 理科基础物理部分

\begin{enumerate}
\item
%% number: 1
%% paperName: 2008年广东理科基础第 1 题 2 分
%% typeId: 1
%% type: 单选题
%% chapter: 电场
%% point: 电场线
%% method: 无
%% score: 2
%% degree: 950
%% duplicateId: 0
%% body:
最早提出用电场线描述电场的物理学家是\xzanswer{C}

\fourchoices[answer=C]
{牛顿}
{伽利略}
{法拉第}
{阿基米德}

%% answer: C
%% memo:
\memoanswer{属于识记的内容。若没有记清楚，也可以采用排除法，除了法拉第以外，其他三个选项提到的物理学家都没有研究电场。}

\item
%% number: 2
%% paperName: 2008年广东理科基础第 2 题 2 分
%% typeId: 1
%% type: 单选题
%% chapter: 力物体的平衡
%% point: 受力分析
%% method: 无
%% score: 2
%% degree: 850
%% duplicateId: 0
%% body:
人站在自动扶梯的水平踏板上，随扶梯斜向上匀速运动，如图所示，以下说法正确的是\xzanswer{A}

\onepicture[w=3.5cm]{figs/02.png}

\fourchoices[answer=A]
{人受到重力和支持力的作用}
{人受到重力、支持力和摩擦力的作用}
{人受到的合外力不为零}
{人受到的合外力方向与速度方向相同}

%% answer: A
%% memo:
\memoanswer{由于人做匀速直线运动，处于平衡状态，人受到的合外力为零，排除了 C、D；只能选 A，才能满足合外力为零这个平衡状态的条件。}

\item
%% number: 3
%% paperName: 2008年广东理科基础第 3 题 2 分
%% typeId: 1
%% type: 单选题
%% chapter: 电路
%% point: 电容
%% method: 无
%% score: 2
%% degree: 950
%% duplicateId: 0
%% body:
关于电容器的电容 $C$、电压 $U$ 和所带电荷量 $Q$ 之间的关系，以下说法正确的是\xzanswer{C}

\fourchoices[answer=C]
{$C$ 由 $U$ 确定}
{$C$ 由 $Q$ 确定}
{$C$ 一定时，$Q$ 与 $U$ 成正比}
{$C$ 一定时，$Q$ 与 $U$ 成反比}

%% answer: C
%% memo:
\memoanswer{$C=\frac{Q}{U}$ 是电容器电容的比值定义式，电容 $C$ 由电容器本身的因素决定，与加在电容器两极板间的电压 $U$ 和极板所带电荷量 $Q$ 都无关，排除 A、B；将公式变形可得 $Q=CU$，可见当 $C$ 一定时，$Q$ 与 $U$ 成正比，选择 C 选项。}

\item
%% number: 4
%% paperName: 2008年广东理科基础第 4 题 2 分
%% typeId: 1
%% type: 单选题
%% chapter: 曲线运动
%% point: 平抛运动
%% method: 无
%% score: 2
%% degree: 850
%% duplicateId: 0
%% body:
水平匀速飞行的直升机上向外自由释放一个物体，不计空气阻力，在物体下落过程中，下列说法正确的是\xzanswer{C}

\fourchoices[answer=C]
{从飞机上看，物体静止}
{从飞机上看，物体始终在飞机的后方}
{从地面上看，物体做平抛运动}
{从地面上看，物体做自由落体运动}

%% answer: C
%% memo:
\memoanswer{物体下落过程中在水平方向上没有受力，速度与直升机的速度相同；在竖直方向上，物体受到重力，做自由落体运动。因此，从飞机上看，物体始终在飞机的正下方，且相对飞机做自由落体运动；从地面上看，物体做平抛运动。所以选择 C 选项。}

\item
%% number: 5
%% paperName: 2008年广东理科基础第 5 题 2 分
%% typeId: 1
%% type: 单选题
%% chapter: 曲线运动
%% point: 万有引力定律
%% method: 无
%% score: 2
%% degree: 850
%% duplicateId: 0
%% body:
人造卫星绕地球做匀速圆周运动，卫星所受万有引力 $F$ 与轨道半径 $r$ 的关系是\xzanswer{D}

\fourchoices[answer=D]
{$F$ 与 $r$ 成正比}
{$F$ 与 $r$ 成反比}
{$F$ 与 $r^{2}$ 成正比}
{$F$ 与 $r^{2}$ 成反比}

%% answer: D
%% memo:
\memoanswer{由万有引力定律可知，卫星所受万有引力 $F=G\frac{Mm}{r^{2}}$，$M$ 是地球的质量，$m$ 是人造卫星的质量，在运动过程中都不变，则 $F$ 与 $r^{2}$ 成反比，选 D。}

\item
%% number: 6
%% paperName: 2008年广东理科基础第 6 题 2 分
%% typeId: 1
%% type: 单选题
%% chapter: 力物体的平衡
%% point: 三力平衡
%% method: 无
%% score: 2
%% degree: 750
%% duplicateId: 0
%% body:
如图所示，质量为 $m$ 的物体悬挂在轻质支架上，斜梁 OB 与竖直方向的夹角为 $\theta$。设水平横梁 OA 和斜梁 OB 作用于 O 点的弹力分别为 $F_{1}$ 和 $F_{2}$，以下结果正确的是\xzanswer{D}

\onepicture[w=3.5cm]{figs/06.png}

\fourchoices[answer=D]
{$F_{1}=mg\sin\theta$}
{$F_{1}=\frac{mg}{\sin\theta}$}
{$F_{2}=mg\cos\theta$}
{$F_{2}=\frac{mg}{\cos\theta}$}

%% answer: D
%% memo:
\memoanswer{以 O 点为研究对象，绳子对 O 点的拉力竖直向下、大小等于 $mg$；水平横梁 OA 对 O 点的弹力 $F_{1}$ 沿水平方向向左，斜梁 OB 对 O 点的弹力 $F_{2}$ 沿 OB 方向向上，如图所示。由平衡条件可得 $F_{1}=mg\tan\theta$，$F_{2}=\frac{mg}{\cos\theta}$，只有 D 选项正确。

\onepicture[w=4.0cm]{figs/06_an.jpg}
}

\item
%% number: 7
%% paperName: 2008年广东理科基础第 7 题 2 分
%% typeId: 1
%% type: 单选题
%% chapter: 曲线运动
%% point: 圆周运动的应用
%% method: 无
%% score: 2
%% degree: 750
%% duplicateId: 0
%% body:
汽车甲和汽车乙质量相等，以相等的速率沿同一水平弯道做匀速圆周运动，甲车在乙车外侧，两车沿半径方向受到的摩擦力分别为 $f_{\text{甲}}$ 和 $f_{\text{乙}}$。以下说法正确的是\xzanswer{A}

\fourchoices[answer=A]
{$f_{\text{甲}}$ 小于 $f_{\text{乙}}$}
{$f_{\text{甲}}$ 等于 $f_{\text{乙}}$}
{$f_{\text{甲}}$ 大于 $f_{\text{乙}}$}
{$f_{\text{甲}}$ 和 $f_{\text{乙}}$ 大小均与汽车速率无关}

%% answer: A
%% memo:
\memoanswer{汽车在水平弯道做匀速圆周运动时，由沿半径方向受到的摩擦力提供向心力 $f=m\frac{v^{2}}{r}$。由于两车质量相等、速率也相等，所需的向心力由轨道半径 $r$ 决定，$r$ 越小所需向心力越大。甲车在乙车外侧，说明 $r_{\text{甲}}>r_{\text{乙}}$，所以 $f_{\text{甲}}<f_{\text{乙}}$，A 正确。}

\item
%% number: 8
%% paperName: 2008年广东理科基础第 8 题 2 分
%% typeId: 1
%% type: 单选题
%% chapter: 曲线运动
%% point: 圆周运动的应用
%% method: 无
%% score: 2
%% degree: 650
%% duplicateId: 0
%% body:
由于地球自转，使得静止在地面的物体绕地轴做匀速圆周运动，对于这些做匀速圆周运动的物体，以下说法正确的是\xzanswer{D}

\fourchoices[answer=D]
{向心力都指向地心}
{速度等于第一宇宙速度}
{加速度等于重力加速度}
{周期与地球自转的周期相等}

%% answer: D
%% memo:
\memoanswer{静止在地面的物体与地球一起转动，转动的周期与地球自转的周期相等，D 正确；物体做圆周运动的圆心在地轴上，向心力一定垂直指向地轴，只有在赤道上的物体向心力才指向地心，A 不正确；第一宇宙速度是环绕速度，比静止在地面上的物体自转的线速度大得多，B 不正确；物体自转的加速度与重力加速度不同，且一般比重力加速度小得多，C 错误。}

\item
%% number: 9
%% paperName: 2008年广东理科基础第 9 题 2 分
%% typeId: 1
%% type: 单选题
%% chapter: 力物体的平衡
%% point: 弹力
%% method: 无
%% score: 2
%% degree: 750
%% duplicateId: 0
%% body:
探究弹力和弹簧伸长的关系时，在弹性限度内，悬挂 $15\UN$ 重物时，弹簧长度为 $0.16\Um$；悬挂 $20\UN$ 重物时，弹簧长度为 $0.18\Um$，则弹簧的原长 $L_{0}$ 和劲度系数 $k$ 分别为\xzanswer{D}

\fourchoices[answer=D]
{$L_{0}=0.02\Um$，$k=500\UNm$}
{$L_{0}=0.10\Um$，$k=500\UNm$}
{$L_{0}=0.02\Um$，$k=250\UNm$}
{$L_{0}=0.10\Um$，$k=250\UNm$}

%% answer: D
%% memo:
\memoanswer{悬挂 $15\UN$ 重物时，$15=k(0.16-L_{0})$……①；悬挂 $20\UN$ 重物时，$20=k(0.18-L_{0})$……②。解①②得 $L_{0}=0.10\Um$，$k=250\UNm$，选择 D。}

\item
%% number: 10
%% paperName: 2008年广东理科基础第 10 题 2 分
%% typeId: 1
%% type: 单选题
%% chapter: 直线运动
%% point: 运动图象
%% method: 无
%% score: 2
%% degree: 850
%% duplicateId: 0
%% body:
如图是某物体做直线运动的 $v-t$ 图象，由图象可得到的正确结果是\xzanswer{B}

\onepicture[w=4.5cm]{figs/10.png}

\fourchoices[answer=B]
{$t=1\Us$ 时物体的加速度大小为 $1.0\Umsq$}
{$t=5\Us$ 时物体的加速度大小为 $0.75\Umsq$}
{第 $3\Us$ 内物体的位移为 $1.5\Um$}
{物体在加速过程的位移比减速过程的位移大}

%% answer: B
%% memo:
\memoanswer{在 $v-t$ 图象中加速度大小由图线斜率的绝对值判断，$t=1\Us$ 时 $a=\frac{3}{2}\Umsq=1.5\Umsq$，A 错误；$t=5\Us$ 时 $a=\frac{3}{4}\Umsq=0.75\Umsq$，B 正确；位移看面积，第 $3\Us$ 内物体以 $3\Ums$ 匀速运动，位移为 $3\Um$，C 错误；由图象可看出，$0\sim2\Us$ 内物体加速，$3\sim7\Us$ 内减速，减速过程图线与坐标轴围成的面积更大，D 错误。}

\item
%% number: 11
%% paperName: 2008年广东理科基础第 11 题 2 分
%% typeId: 1
%% type: 单选题
%% chapter: 功和能
%% point: 动能定理
%% method: 无
%% score: 2
%% degree: 650
%% duplicateId: 0
%% body:
一个 $25\Ukg$ 的小孩从高度为 $3.0\Um$ 的滑梯顶端由静止开始滑下，滑到底端时的速度为 $2.0\Ums$，取 $g=10\Umsq$，关于力对小孩做的功，以下结果正确的是\xzanswer{A}

\fourchoices[answer=A]
{合外力做功 $50\UJ$}
{阻力做功 $500\UJ$}
{重力做功 $500\UJ$}
{支持力做功 $50\UJ$}

%% answer: A
%% memo:
\memoanswer{根据动能定理，合外力做的功等于动能的变化量，小孩初速度为 $0$、末速度为 $2.0\Ums$，则 $W_{\text{合}}=\frac{1}{2}mv^{2}=\frac{1}{2}\times25\times2.0^{2}\UJ=50\UJ$，A 正确；其余力做的功无法由已知条件判断。}

\item
%% number: 12
%% paperName: 2008年广东理科基础第 12 题 2 分
%% typeId: 1
%% type: 单选题
%% chapter: 运动和力
%% point: 牛顿第二定律
%% method: 无
%% score: 2
%% degree: 750
%% duplicateId: 0
%% body:
质量为 $m$ 的物体从高处静止释放后竖直下落，在某时刻受到的空气阻力为 $f$，加速度为 $a=\frac{1}{3}g$，则 $f$ 的大小是\xzanswer{B}

\fourchoices[answer=B]
{$f=\frac{1}{3}mg$}
{$f=\frac{2}{3}mg$}
{$f=mg$}
{$f=\frac{4}{3}mg$}

%% answer: B
%% memo:
\memoanswer{根据牛顿第二定律 $mg-f=ma$，而 $a=\frac{1}{3}g$，解得 $f=\frac{2}{3}mg$，选 B。}

\item
%% number: 13
%% paperName: 2008年广东理科基础第 13 题 2 分
%% typeId: 1
%% type: 单选题
%% chapter: 电路
%% point: 电路中的图象
%% method: 无
%% score: 2
%% degree: 850
%% duplicateId: 0
%% body:
在“测定电源电动势和内阻”的实验中，某同学根据实验数据，作出了正确的 $U-I$ 图象。如图所示，其中图线斜率绝对值的物理含义是\xzanswer{B}

\onepicture[w=4.5cm]{figs/13.png}

\fourchoices[answer=B]
{短路电流}
{电源内阻}
{电源电动势}
{全电路电阻}

%% answer: B
%% memo:
\memoanswer{$U-I$ 图象表示路端电压随外电路总电流变化的关系，图线斜率绝对值表示电源内阻，B 正确；纵轴截距表示电源电动势，横轴截距表示短路电流。}

\item
%% number: 14
%% paperName: 2008年广东理科基础第 14 题 2 分
%% typeId: 1
%% type: 单选题
%% chapter: 电路
%% point: 伏安法测电阻
%% method: 无
%% score: 2
%% degree: 850
%% duplicateId: 0
%% body:
某同学用伏安法测小灯泡的电阻时，误将电流表和电压表接成如图所示的电路。接通电源后，可能出现的情况是\xzanswer{D}

\onepicture[w=4.0cm]{figs/14.png}

\fourchoices[answer=D]
{电流表烧坏}
{电压表烧坏}
{小灯泡烧坏}
{小灯泡不亮}

%% answer: D
%% memo:
\memoanswer{由于电压表内阻一般较大、电流表内阻一般较小，根据串联分压和并联分流的规律，最终流过小灯泡的电流很小，故小灯泡不亮，D 正确。}

\item
%% number: 15
%% paperName: 2008年广东理科基础第 15 题 2 分
%% typeId: 1
%% type: 单选题
%% chapter: 电路
%% point: 电阻定律
%% method: 无
%% score: 2
%% degree: 950
%% duplicateId: 0
%% body:
关于电阻率的说法正确的是\xzanswer{B}

\fourchoices[answer=B]
{电阻率与导体的长度有关}
{电阻率与导体的材料有关}
{电阻率与导体的形状有关}
{电阻率与导体的横截面积有关}

%% answer: B
%% memo:
\memoanswer{电阻率是描述导体导电性能好坏的物理量，仅与导体的材料和温度等因素有关，与导体的长度、形状、横截面积无关，B 正确。}

\item
%% number: 16
%% paperName: 2008年广东理科基础第 16 题 2 分
%% typeId: 1
%% type: 单选题
%% chapter: 电场
%% point: 带电粒子的偏转
%% method: 无
%% score: 2
%% degree: 650
%% duplicateId: 0
%% body:
空间存在竖直向上的匀强电场，质量为 $m$ 的带正电的微粒水平射入电场中，微粒的运动轨迹如图所示。在相等的时间间隔内\xzanswer{C}

\onepicture[w=3.5cm]{figs/16.png}

\fourchoices[answer=C]
{重力做的功相等}
{电场力做的功相等}
{电场力做的功大于重力做的功}
{电场力做的功小于重力做的功}

%% answer: C
%% memo:
\memoanswer{带正电的微粒在竖直方向上受到竖直向下的重力和竖直向上的电场力，由运动轨迹向上偏可知电场力大于重力，因此 C 正确；在相等的时间间隔内微粒在竖直方向上的位移不同，A、B 均错误。}

\item
%% number: 17
%% paperName: 2008年广东理科基础第 17 题 2 分
%% typeId: 1
%% type: 单选题
%% chapter: 磁场
%% point: 洛伦兹力
%% method: 无
%% score: 2
%% degree: 850
%% duplicateId: 0
%% body:
有关洛伦兹力和安培力的描述，正确的是\xzanswer{B}

\fourchoices[answer=B]
{通电直导线在匀强磁场中一定受到安培力的作用}
{安培力是大量运动电荷所受洛伦兹力的宏观表现}
{带电粒子在匀强磁场中运动受到的洛伦兹力做正功}
{通电直导线在磁场中受到的安培力方向与磁场方向平行}

%% answer: B
%% memo:
\memoanswer{当通电直导线放置方向与匀强磁场方向在同一直线上时，不受安培力作用，A 错误；安培力可看成导体内大量定向移动的电荷共同受到洛伦兹力的宏观表现，B 正确；洛伦兹力始终与运动方向垂直，不做功，C 错误；由左手定则可知，安培力方向与磁场方向垂直，D 错误。}

\item
%% number: 18
%% paperName: 2008年广东理科基础第 18 题 2 分
%% typeId: 1
%% type: 单选题
%% chapter: 磁场
%% point: 带电粒子在磁场中的运动
%% method: 无
%% score: 2
%% degree: 750
%% duplicateId: 0
%% body:
电子在匀强磁场中做匀速圆周运动，下列说法正确的是\xzanswer{D}

\fourchoices[answer=D]
{速率越大，周期越大}
{速率越小，周期越大}
{速度方向与磁场方向平行}
{速度方向与磁场方向垂直}

%% answer: D
%% memo:
\memoanswer{电子在匀强磁场中运动，由洛伦兹力提供向心力 $evB=m\frac{v^{2}}{r}$，运动周期 $T=\frac{2\pi r}{v}=\frac{2\pi m}{eB}$，与速率无关，A、B 均错误；由左手定则可知，电子的速度方向与磁场方向垂直，D 正确，C 错误。}

\end{enumerate}

\end{document}
